\subsection{A complete field-affine proof compiler}\label{sec:realize:compiler}

\begin{definition}[Bit-committed field variables]\label{def:realize:field-commit}
To register a hidden field element $x\in L$, commit
independently to its $n$ fixed-basis coordinates
\[
 C_i(x):=\mathsf{Com}(d_i(x);r_i),\qquad
 0\leq i<n,\quad r_i\overset{\$}{\leftarrow}\mathbb Z_P.
\]
Write $\mathbf C(x)$ for the $n$-tuple.
The $n$ bitness OR proofs establish that every
registered group element has a known binary opening.
Registration occurs before the outer random challenge
whose soundness argument relies on the field value.
\end{definition}

\begin{lemma}[Registered coordinates reconstruct one field element]\label{lem:realize:field-reconstruct}
Given extracted binary openings $(b_i,r_i)$ of
every $C_i(x)$, there is a unique associated field value
\[
 x^{\rm ext}:=\sum_{i=0}^{n-1}\gamma_i b_i\in L.
\]
\end{lemma}
\begin{proof}
The fixed basis coefficients $b_i\in\F_2$ determine
the displayed field element. Uniqueness follows
from the uniqueness of basis expansions;
no second field element can have the same
coordinate vector.
\end{proof}

\begin{definition}[Affine proof compiler]\label{def:realize:affine-compiler}
For registered hidden field variables $x_1,\ldots,x_s$
and a public affine equation $\sum_ja_jx_j=b$,
the compiler:
\begin{enumerate}[label=\textup{(\roman*)}]
\item computes all public binary coefficients
$M_{k,j,i}=d_k(a_j\gamma_i)$ and $b_k=d_k(b)$;
\item for each row $k$, selects the set
$T_k=\{(j,i):M_{k,j,i}=1\}$;
\item forms the public group product
$C_{T_k}=\prod_{(j,i)\in T_k}C_i(x_j)$;
\item proves the integer opening of $C_{T_k}$
lies in $\mathcal S_{|T_k|,b_k}$
using the OR protocol of
Definition~\ref{def:realize:or}.
\end{enumerate}
Each registered bit must additionally have a valid
bitness proof from Definition~\ref{def:realize:parity-proof}.
No secret field element is transmitted or opened.
\end{definition}

\begin{lemma}[All affine row statements are publicly computable]\label{lem:realize:rows-public}
Given the public group commitments, public
basis, coefficients $(a_j)_j$ and right side $b$,
a verifier can compute every $C_{T_k}$ and
candidate set $\mathcal S_{|T_k|,b_k}$
without knowing a hidden field element.
\end{lemma}
\begin{proof}
Every $M_{k,j,i}$ equals the publicly
computable coordinate $d_k(a_j\gamma_i)$.
The set $T_k$ and $b_k=d_k(b)$ therefore
depend only on public data. Its aggregate
commitment is the product of public
commitments indexed by $T_k$, and the
candidate set depends only on the publicly
known row cardinality and target bit.
\end{proof}

\begin{theorem}[Perfect completeness of the field-affine compiler]\label{thm:realize:compiler-complete}
If every registered commitment opens to the correct
coordinate of $x_j$ and
$\sum_ja_jx_j=b$ holds in $L$, then
all bitness and compiled parity proofs
accept for any public coefficient choices.
\end{theorem}
\begin{proof}
Every bitness proof has a known binary opening
and is complete by
Lemma~\ref{lem:realize:or-complete}.
By Theorem~\ref{thm:realize:affine-to-parity},
each of the $n$ compiled parity rows
is true for these coordinate bits.
Theorem~\ref{thm:realize:parity-complete}
makes every corresponding parity
OR proof complete. All checks pass.
\end{proof}

\begin{theorem}[Extracted field-affine correctness]\label{thm:realize:compiler-extract}
Suppose all registered commitments have
extracted binary openings, and each compiled
parity OR proof has an extracted opening
to one of its allowed candidates.
If the public field equation is false for
the extracted field elements
$x_j^{\rm ext}$, then the extracted openings
give a polynomial-time solution to the
setup discrete-logarithm problem.
\end{theorem}
\begin{proof}
If the field equation is false,
Theorem~\ref{thm:realize:affine-to-parity}
implies that at least one output-coordinate
parity row is false for the extracted bits.
Apply
Theorem~\ref{thm:realize:parity-binding}
to that row's extracted bit openings
and allowed aggregate opening.
It computes $\log_g h$ in polynomial time,
contradicting setup hardness.
\end{proof}

\begin{lemma}[Faithful compilation of the Chapter 9 checker]\label{lem:realize:chapter9-affine}
Each hidden-field equation in
Definition~\ref{def:zkrs:affine-checks}
can be expressed in the affine input
format of Definition~\ref{def:realize:affine-compiler}
after its stipulated public messages
and challenges are fixed.
The trace constraints are public
and require no hidden-field proof.
\end{lemma}
\begin{proof}
Lemma~\ref{lem:zkrs:claim-in-affine}
handles the claimed-value and
incoming equations.
Lemma~\ref{lem:zkrs:out-affine}
handles the outgoing equation, and
Lemma~\ref{lem:zkrs:anchor-affine}
handles both anchored equations.
In each case, all coefficients multiplying
hidden field variables belong to the
publicly known field $L$ once the
respective public messages and
challenges have been received.
Lemma~\ref{lem:zkrs:public-traces}
makes every trace-zero test directly
computable from public values.
\end{proof}

\begin{lemma}[Packed-column equality is also field affine in bits]\label{lem:realize:packing-link}
For a binary witness column $(w_{v,j})_{v\in I}$,
its packed field value
$\widehat w_j=\sum_v\beta_vw_{v,j}$
is equivalent to $n$ binary linear
relations between the bits of $\widehat w_j$
and the witness bits $w_{v,j}$.
\end{lemma}
\begin{proof}
Apply each coordinate map $d_k$ to
$\widehat w_j=\sum_v\beta_vw_{v,j}$.
Since every $w_{v,j}\in\F_2$,
linearity gives
\[
 d_k(\widehat w_j)
   =\sum_v d_k(\beta_v)w_{v,j}
   =w_{k,j},
\]
using $d_k(\beta_v)=\delta_{kv}$.
Thus the $n$ equalities are simply
$ d_k(\widehat w_j)=w_{k,j}$.
Conversely these $n$ coordinate
equalities imply the field equality
by unique basis reconstruction.
\end{proof}

\begin{remark}[Separate global witness binding]
The last lemma allows commitment-level linking
between the packed table and a separately
committed binary table if both are exposed to
the concrete bit-proof system. It does not
automatically link those commitments to an
\emph{external KBFold polynomial oracle}
whose authentication structure is different.
That link needs an explicitly specified
cross-commitment statement and a proof.
\end{remark}
